\section{Design Patterns}

Come richiesto dalle specifiche progettuali d'esame, la presente sezione descrive l'applicazione di due \textit{Design Pattern} (GoF --- Gang of Four) al fine di risolvere criticità architetturali del modello a oggetti, garantendo manutenibilità, basso accoppiamento ed elevata coesione.

% -----------------------------------------------------------
\subsection{Design Pattern 1: Observer Pattern}

\subsubsection{Motivazione e Criticità di Design}
Nel sistema \textbf{MyAma}, quando una prenotazione subisce una transizione di stato (ad esempio da \textit{In attesa} a \textit{Confermata}, oppure viene \textit{Annullata} o \textit{Completata} dall'autista), molteplici sottosistemi indipendenti devono reagire tempestivamente:
\begin{itemize}
    \item Il modulo di notifica al cittadino (email / SMS / push).
    \item La dashboard dell'autista (per aggiornare l'itinerario e il carico residuo).
    \item Il registro operativo della sede.
\end{itemize}
L'accoppiamento diretto della classe \texttt{Prenotazione} con tutte le classi concrete di notifica e aggiornamento violerebbe il principio \textit{Open/Closed} e aumenterebbe l'accoppiamento.

\subsubsection{Soluzione Progettuale}
L'applicazione del pattern \textbf{Observer} permette di definire una dipendenza uno-a-molti: la classe \texttt{Prenotazione} funge da \textit{Subject} (o \textit{Observable}), notificando tutti gli \textit{Observer} registrati (\texttt{EmailNotifier}, \texttt{DashboardNotifier}, \texttt{LogNotifier}) ogni qualvolta il proprio stato muta, mantenendo totale disaccoppiamento.

\begin{figure}[H]
    \centering
    % \includegraphics[width=0.85\textwidth]{figure/pattern_observer.png}
    \fbox{\parbox[c][6.5cm][c]{0.8\textwidth}{\centering \textsf{[ Inserire qui il Diagramma UML dell'Observer Pattern applicato a MyAma esportato da Visual Paradigm ]}}}
    \caption{Struttura UML del Design Pattern Observer applicato a MyAma}
    \label{fig:pattern_observer}
\end{figure}

% -----------------------------------------------------------
\subsection{Design Pattern 2: Strategy Pattern}

\subsubsection{Motivazione e Criticità di Design}
L'assegnazione ottimale dei veicoli e la pianificazione degli slot di ritiro per i rifiuti ingombranti possono seguire diversi algoritmi operativi a seconda delle condizioni aziendali:
\begin{itemize}
    \item \textbf{Strategia per Minimizzazione del Volume/Peso:} Ottimizza il riempimento del mezzo per massimizzare i ritiri per singolo viaggio.
    \item \textbf{Strategia per Prossimità Territoriale / Distanza:} Minimizza i chilometri percorsi dall'autista all'interno del medesimo CAP.
    \item \textbf{Strategia Bilanciata per Carico di Lavoro:} Distribuisce equamente le prenotazioni tra tutti gli autisti in turno.
\end{itemize}
Codificare questi algoritmi direttamente all'interno della classe di controllo (\texttt{GestoreAssegnazione}) comporterebbe complessi blocchi condizionali (\texttt{if-else}/\texttt{switch}) difficilmente manutenibili ed estendibili.

\subsubsection{Soluzione Progettuale}
L'applicazione dello \textbf{Strategy Pattern} incapsula ciascun algoritmo di allocazione in classi separate conformi all'interfaccia comune \texttt{AssegnazioneStrategy}. In questo modo il contesto può variare dinamicamente la strategia di scheduling a tempo di esecuzione senza modificare il codice client.

\begin{figure}[H]
    \centering
    % \includegraphics[width=0.85\textwidth]{figure/pattern_strategy.png}
    \fbox{\parbox[c][6.5cm][c]{0.8\textwidth}{\centering \textsf{[ Inserire qui il Diagramma UML dello Strategy Pattern applicato a MyAma esportato da Visual Paradigm ]}}}
    \caption{Struttura UML del Design Pattern Strategy per l'assegnazione risorse}
    \label{fig:pattern_strategy}
\end{figure}
